\documentclass[11pt]{article}
\usepackage[letterpaper,margin=1in]{geometry}
\usepackage[T1]{fontenc}
\usepackage{lmodern,microtype}
\usepackage{amsmath,amssymb,amsthm,booktabs,array}
\usepackage{siunitx}
\usepackage{xurl}
\usepackage[colorlinks=true,linkcolor=blue!50!black,citecolor=blue!50!black,urlcolor=blue!50!black]{hyperref}
\usepackage{xcolor}
\sisetup{group-separator={,},group-minimum-digits=4}
\newtheorem{theorem}{Theorem}[section]
\newtheorem{proposition}[theorem]{Proposition}
\newtheorem{lemma}[theorem]{Lemma}
\newtheorem{corollary}[theorem]{Corollary}
\newcommand{\F}{F_3}
\newcommand{\Slice}{\mathcal S}
\newcommand{\Chart}{\mathcal H}
\newcommand{\word}[1]{\texttt{#1}}
\DeclareMathOperator{\rot}{rot}
\newcommand{\lam}{\frac{926683}{1814400}}
\newcommand{\gam}{\frac{30571}{1814400}}
\title{A smaller upper bound for superpermutations}
\author{Jay Pantone}
\date{October 2026}
\hypersetup{pdftitle={A smaller upper bound for superpermutations},pdfauthor={Jay Pantone}}

\begin{document}
\maketitle

\begin{abstract}
We construct superpermutations of length
\[
 n!+(n-1)!+(n-2)!+
 \left(\frac{926683}{1814400}+o(1)\right)(n-3)!.
\]
The coefficient is approximately $0.510738$, improving our earlier
$43/80$. The construction uses the same completion and transport operations,
with a different finite selection and a smaller family of connector cycles.
An explicit version applies for every $n\ge14$. The finite certificate,
its extension to all larger alphabets, and the resulting word-existence
bounds are verified in Lean. We also give explicit words on ten through
thirteen symbols, checked directly for permutation coverage.
\end{abstract}

\paragraph{AI disclosure.}
The research presented here, the Lean formalization, and this document were generated entirely by AI systems under Jay Pantone's direction. The inherited mathematical work and software retain their original credits. Independent AI review is not human peer review.

\section{The bound}

A word can contain many permutations without repeating all of their
letters. For example, \word{123121321} contains each of the six permutations
of $\{1,2,3\}$ as a contiguous substring. A word with this property for an
alphabet of $n$ symbols is a \emph{superpermutation}. We write $S(n)$ for
the least possible length and put
\[
 \F(n)=n!+(n-1)!+(n-2)!.
\]
Our concern here is the next factorial term in an upper bound for $S(n)$.

The construction of Pantone~\cite{pantone} gives coefficient $43/80$ on
$(n-3)!$. Its main idea is to cover most permutations with long pieces
whose ends fit together, repair the omitted permutations, and then join
the resulting closed trails. There are two substantial costs. Omitting
permutations from a piece makes it easier to join, but those permutations
must later be restored. Connector cycles join many of the remaining
trails at once. The coefficient is the sum of these two costs.

Halfar~\cite{halfar} improved the finite selection used in this construction
and introduced replacements inside prescribed cyclic blocks. Our new
selection uses that replacement method, followed by an optimized family
of connector cycles. The completion, transport, and final joining
arguments are those of the earlier construction. The improvement is in
the finite object to which they are applied.

\begin{theorem}\label{thm-asymptotic}
For every $\varepsilon>0$ there is an $N$ such that
\[
 S(n)\le \F(n)+\left(\lam+\varepsilon\right)(n-3)!
 \qquad\text{for every }n\ge N.
\]
\end{theorem}

The reduction from $43/80$ is exactly $48557/1814400$. Halfar's
October 2026 construction has coefficient $1771/3456$; the reduction from
that coefficient is $773/453600$. The coefficient proved here is still larger
than $1/2$. The comparisons in this paper concern the specified
constructions and announced results; no current-record claim is made.

The finite form of the theorem is useful independently of its asymptotic
consequence. The parameter $\ell$ is the number of letters overlapped when
we join the final pieces.

\begin{theorem}\label{thm-finite}
Let $n\ge14$, let $0\le\ell\le n-4$, and put
$c=\gam(n-4)!$. Then
\[
 \begin{split}
 S(n)\le {}&\F(n)+\lam(n-3)!+(n-4-\ell)c\\
 &\hspace{20mm}+\ell\min\left(c,\frac{(n-1)!}{(n-1-\ell)!}\right).
 \end{split}
\]
\end{theorem}

Taking $\ell=n-11$ gives the values in Table~\ref{tab-finite}. These are
bounds on the lengths of words whose existence is proved by the
construction. There is no need to write out those words to prove the
bounds. The much smaller files described in Section~\ref{sec-words}
contain actual words on ten through thirteen symbols.

\begin{table}[ht]
\centering
\begin{tabular}{rrr}
\toprule
$n$ & $\ell$ & Upper bound on $S(n)$\\
\midrule
14 & 3 & \num{93905133768}\\
15 & 4 & \num{1401329128342}\\
16 & 5 & \num{22320881220264}\\
17 & 6 & \num{377963186649048}\\
18 & 7 & \num{6779662770937056}\\
19 & 8 & \num{128414015952702480}\\
20 & 9 & \num{2561133914781361920}\\
\bottomrule
\end{tabular}
\caption{Consequences of Theorem~\ref{thm-finite}, all verified in Lean.}
\label{tab-finite}
\end{table}

We first describe the pieces and the operations that allow a finite
construction to work on every larger alphabet. The new selection and
its connector certificate are then small enough to describe in terms of
five source trails. The final section treats the explicit words and their
verification.

\section{Pieces that fit together}\label{sec-pieces}

Fix an alphabet of $k\ge4$ symbols. For a permutation $p$, let $[p]$ be its
\emph{cyclic class}, the set of its $k$ rotations, and let $\rot$ denote
left rotation. Consecutive rotations overlap in $k-1$ letters, so all of
$[p]$ can be read in a word of length $2k-1$.

Choose a distinguished symbol $s$, and let $x$ order the other $k-1$
symbols. The \emph{2-loop} determined by $(x,s)$ visits
\[
 [xs],\ [\rot(x)s],\ \ldots,\ [\rot^{k-2}(x)s]
\]
in that order. Within a class it starts at $\rot^j(x)s$ and follows
rotations until it reaches $s\rot^j(x)$. The transition to
$\rot^{j+1}(x)s$ then costs two new letters. Replacing $x$ by a rotation
changes only the starting point of the loop. With $s$ fixed, the
$(k-2)!$ such loops partition the permutations.

A \emph{slice} $\Slice(x,s;L)$ takes the first $L$ classes of this loop.
Its word has length
\[
 2k-1+(L-1)(k+1)=(k+1)L+k-2.
\]
For example, $\Slice(123,4;1)$ spells \word{1234123}. It contains the
class $[1234]$ and omits the classes $[2314]$ and $[3124]$. A slice is
\emph{full} when $L=k-1$ and \emph{short} when $L=k-3$.

The first and last $k-3$ letters of a slice are its \emph{head} and
\emph{tail}. A slice can be followed by another whose head is its tail,
with an overlap of $k-3$ letters. In the example the head is \word{1}
and the tail is \word{3}. We regard a slice as a directed edge from its
head to its tail.

A \emph{selection} chooses a full or short slice from every 2-loop with
distinguished symbol $s$. Its \emph{charge} is twice the number of short
slices. Thus the charge $Q$ counts the cyclic classes that the selection
omits. Our selections have balanced endpoint graphs, so their slices can
be partitioned into closed trails. We retain a particular such partition,
since the way the trails behave when a letter is added is essential.
For a closed trail $T$, write $f(T)$ and $t(T)$ for its numbers of full
and short slices, and call $f(T)-t(T)$ its \emph{signed excess}.

One simple closed trail will be used repeatedly. Let $y$ be a word of
$h$ distinct symbols, and let $a,s$ be two further symbols. The full slices
\[
 \Slice(\rot^j(y)a,s;h+1),\qquad 0\le j<h,
\]
form a closed trail, which we denote by $\Chart(y;a,s)$. For example,
$\Chart(123;4,5)$ has heads \word{12}, \word{23}, and \word{31}, and
its tails are the same three words in the next cyclic order. We call this
a \emph{full chart}. Its charge is zero and its signed excess is $h$.

\subsection{Completion}

Completion adds a final symbol $z$ and fills in every omitted cyclic
class. If $x=x_0\cdots x_{k-2}$, write $x^{(i)}$ for the result of
inserting $z$ immediately before $x_i$. The slice $\Slice(x,s;L)$ gives
$k-1$ inserted slices. The one with base $x^{(i)}$ has $L+1$ classes if
$i<L$, and $L$ classes otherwise. For each omitted class
$[\rot^i(x)s]$, add the full slice with distinguished symbol $z$ and base
$\rot^i(x)s$.

Deleting $z$ explains the construction. Each old cyclic class has $k$
extensions, according to the position of $z$ in its cyclic order. In an
inserted 2-loop, deletion lists the old classes in order, except that one
class occurs twice. The extra class in an inserted slice accounts for this
repetition precisely when the repeated class was selected. A repair slice
contains all $k$ extensions of one omitted class. Consequently completion
contains every new cyclic class exactly once.

The inserted and repair slices from one old slice form $k-2$ paths.
Their initial vertices are the ways to insert $z$ into the old head, and
their terminal vertices are the ways to insert $z$ into the old tail.
Number the positions $0,\ldots,k-3$. A short slice sends position $p$ to
$p+1$, and a full slice sends it to $p-1$, modulo $k-2$. Thus the paths
along a closed trail $T$ close after following an orbit of the translation
\[
 p\longmapsto p-f(T)+t(T)\pmod{k-2}.
\]
In particular, completion produces $\gcd(k-2,f(T)-t(T))$ closed trails
from $T$.

If the resulting alphabet has size $n=k+1$, completion has
$(n-2)!+Q$ slices, covering $(n-1)!$ cyclic classes. After its head, a
slice with $L$ classes contributes $(n+1)L+1$ letters. The total
contribution of the completed slices, before any cost for opening or
joining trails, is therefore
\begin{equation}\label{eq-slice-cost}
 (n+1)(n-1)!+(n-2)!+Q=\F(n)+Q.
\end{equation}
This identity is the reason for measuring a selection by its charge.

\subsection{Transport}

To use the same selection on a larger alphabet, we need an operation
that again produces full and short slices. \emph{Transport} adds a new
ordinary symbol $w$ while retaining $s$. A full slice gives all $k-1$
full slices obtained by inserting $w$ before a symbol of $x$. For a short
slice, write $x=uab$. Use the corresponding short slices, except that
the base $uawb$ is replaced by its rotation $buaw$.

For example, transporting $\Slice(123,4;1)$ with the new symbol $5$
gives short slices with bases \word{5123}, \word{1523}, and \word{3125}.
Their endpoint paths are
\[
 \word{51}\longrightarrow\word{35},\qquad
 \word{15}\longrightarrow\word{31}\longrightarrow\word{53}.
\]
The same insertion-position rule holds as for completion. The exceptional
rotation is what makes the second path fit together.

Transport gives one slice in each extension of the old 2-loop, so it
preserves the selection property. It multiplies charge by $k-1$.
If $k-2$ divides the signed excess of each old trail, each trail gives
$k-2$ closed trails, one for each insertion position. Each new trail has
signed excess
\begin{equation}\label{eq-excess}
 \frac{k-1}{k-2}\bigl(f(T)-t(T)\bigr),
\end{equation}
which is divisible by $k-1$. This condition can therefore be maintained
through arbitrarily many transports. These transport and completion
statements, including the occurrence-by-occurrence inventories, are
proved in Pantone~\cite{pantone}.

\section{Connectors and the cost of joining}\label{sec-join}

After completion to $n$ symbols, the endpoint words have length
$h=n-3$. A \emph{connector cycle} is the cycle of rotations of a word
of $h$ distinct symbols. Each of its $h$ edges appends one letter. If
a connector meets a closed trail at an endpoint, it can be spliced into
that trail. A connector may meet several trails, and this is where the
main saving occurs.

There are two kinds of connector incidences used by the construction.
A connector containing the completion symbol $z$ meets an actual endpoint
of a completed trail. A connector not containing $z$ is the cycle of
$y$ for a parent full chart $\Chart(y;a,s)$. This cycle meets every trail
obtained by completing that chart. Both kinds of incidence persist under
transport. A connector of length $h$ gives $h$ new connector cycles of
length $h+1$, one for each gap in which the new ordinary symbol can be
inserted. The first case follows by tracking the two insertion positions,
for the ordinary symbol and for $z$. The second follows by transporting
the full chart. This is the connector transport theorem of
Pantone~\cite{pantone}; its two incidence hypotheses are part of the
statement.

For clarity, we record exactly the consequence that is needed here.

\begin{proposition}[Transport and joining]\label{prop-transport}
Let $n_0\ge5$. Suppose a balanced selection on $n_0-1$ symbols, with distinguished
symbol $s$, has charge $Q_0$ and a partition into closed trails whose
signed excesses are divisible by $n_0-3$. Suppose its completion is met
by $c_0$ connector cycles of length $n_0-3$, none containing $s$, with
each designated incidence of one of the two kinds described above.
Then for every $n\ge n_0$ the transported selection has charge
\[
 Q_n=Q_0\frac{(n-3)!}{(n_0-3)!},
\]
and its completion has a connector cover of size
\[
 c_n=c_0\frac{(n-4)!}{(n_0-4)!}.
\]
For every $0\le\ell\le n-4$, it gives a superpermutation of length at most
\[
 \F(n)+Q_n+(n-3)c_n+(n-4-\ell)c_n
       +\ell\min\left(c_n,\frac{(n-1)!}{(n-1-\ell)!}\right).
\]
\end{proposition}

\begin{proof}
The formulas for $Q_n$ and $c_n$ follow by multiplying the factors in
successive transports. The divisibility in~\eqref{eq-excess} preserves
closed trails, and the connector transport theorem preserves their cover.
It remains to account for the final joining cost.

At a fixed size $n$, put the completed slices and the connector edges
into a directed multigraph on endpoint words. The graph is balanced.
Every connected component contains a connector, because every completed
trail meets one. Hence the number $b$ of components is at most $c_n$.
Choose a connector in each component, take an Eulerian circuit, and remove
one edge of the chosen connector. The resulting word retains every slice.
By~\eqref{eq-slice-cost}, the sum of the lengths of these $b$ words is
\[
 \F(n)+Q_n+(n-3)c_n+(n-4)b.
\]
Here each word starts with $n-3$ letters, and the removed edge saves one.

There is a choice of the removed edges and an ordering into at most
\[
 J\le\min\left(b,\frac{(n-1)!}{(n-1-\ell)!}\right)
\]
trails of words such that successive words in a trail overlap in $\ell$
letters. This is the balanced-cut lemma in Pantone~\cite{pantone}.
Its application can be seen directly from the available cuts. Regard a
cut as an edge from the first $\ell$ letters of its word to the last
$\ell$ letters. Giving equal weight to all cuts on a chosen connector
balances every vertex. The lemma rounds these weights to one cut per
component with absolute imbalance at most two at each vertex. The chosen
edges then decompose into at most one trail per vertex. All vertices are
words of $\ell$ distinct symbols other than $s$, so the stated vertex
count applies. The case $\ell=0$ is simply concatenation.

The overlaps save $\ell(b-J)$ letters. Thus the total length is at most
\[
 \F(n)+Q_n+(n-3)c_n+(n-4-\ell)b+\ell J.
\]
Since $n-4-\ell\ge0$, substituting $b\le c_n$ gives the result.
\end{proof}

The two quantities to minimize are now apparent. The charge contributes
$Q_0/(n_0-3)!$ to the coefficient, and the connectors contribute
$c_0/(n_0-4)!$. The remaining joining cost is asymptotically smaller.
Our finite construction has $n_0=14$ and
\begin{equation}\label{eq-base-numbers}
 Q_0=\num{19714464},\qquad c_0=\num{61142}.
\end{equation}
The next section establishes the selection, its trail partition, and
the connector cover that give these numbers.

\section{The finite construction}\label{sec-base}

It is useful to separate the ordinary symbols from the two symbols with
special roles. The final selection has ordinary alphabet
$\{0,1,\ldots,11\}$ and distinguished symbol $s=12$. Completion adds
$z=13$. A row of the finite data consists of an ordering $x$ of the
ordinary alphabet and an indication of whether the slice is full or short.
The data specify the order of rows along each source trail.

\subsection{Replacing full charts}

The starting selection, due to Halfar, has eight ordinary symbols.
It consists of two closed trails $H_0,H_1$, together with 48 full charts.
The chart words have length seven. The source data have the counts in
Table~\ref{tab-sources}. In particular, the two trails occupy 4704
2-loops and the full charts occupy the other $48\cdot7=336$ loops.
The total is $7!=5040$.

Each of the 48 charts can be treated independently. Transport such a
chart three times, with three new ordinary symbols. Its 5040 descendant
2-loops are characterized by the cyclic order that remains after deleting
the four symbols added to the chart word. This characterization makes it
possible to replace the slices inside the block without affecting any
other block.

On eleven ordinary symbols we use three closed trails $B_0,B_1,B_2$.
Together they occupy 4074 of the 5040 loops of one such block. The other
966 loops are retained as a list of cyclic bases. On adding the twelfth
ordinary symbol, transport the three trails and fill each of the 966
remaining loops with its full chart. Relabel the original seven symbols
to do this in each of the 48 blocks.

\begin{table}[ht]
\centering
\begin{tabular}{lrrrr}
\toprule
Source trail & Ordinary symbols & Slices & Charge & Signed excess\\
\midrule
$H_0$ & 8 & 2280 & 1104 & 1176\\
$H_1$ & 8 & 2424 & 1248 & 1176\\
\addlinespace
$B_0$ & 11 & 714 & 588 & 126\\
$B_1$ & 11 & 1330 & 560 & 770\\
$B_2$ & 11 & 2030 & 910 & 1120\\
\bottomrule
\end{tabular}
\caption{The five finite source trails. Charge is twice the number of
short slices, so signed excess is the number of slices minus the charge.}
\label{tab-sources}
\end{table}

The finite assertions about these lists are straightforward to state.
Every row uses the correct symbols, every slice is full or short, and
consecutive endpoints agree, including the last and first. The rows of
$H_0,H_1$ and the 48 charts represent all 5040 cyclic bases on eight
ordinary symbols. Similarly, the rows of $B_0,B_1,B_2$ and the 966
remaining bases represent all 5040 cyclic bases of the prescribed block
on eleven symbols. Representatives are compared up to rotation. These
coverage assertions, as well as the counts in Table~\ref{tab-sources},
are checked from the literal lists in the accompanying Lean certificate.

The block replacement preserves a selection, not merely its number of
rows. Indeed, transport takes the cyclic base of a row to precisely its
insertion extensions. Consequently equality of the represented cyclic
bases is preserved by transport and by relabeling. The two finite
coverage checks therefore prove coverage of the full ordinary alphabet
after replacement. There are exactly $11!$ rows, one for each cyclic
base on twelve ordinary symbols, so no base can be repeated.

The resulting row and charge counts can now be obtained without expanding
the full selection. The unchanged part comes from transporting $H_0,H_1$
four times. Each replaced block has
$4074\cdot11+966\cdot11=\num{55440}$ rows and charge
$2058\cdot11=\num{22638}$. Thus
\[
 \begin{aligned}
 \text{rows}&=4704(8\cdot9\cdot10\cdot11)+48(55440)
             =\num{39916800}=11!,\\
 Q_0&=2352(8\cdot9\cdot10\cdot11)+48(22638)
             =\num{19714464}.
 \end{aligned}
\]
This proves the charge in~\eqref{eq-base-numbers}.

\subsection{The closed trails}

The two Halfar trails have signed excess 1176, divisible by seven.
Their four transports have $7,8,9,10$ insertion positions, respectively,
so they give
\[
 2\cdot7\cdot8\cdot9\cdot10=\num{10080}
\]
closed trails on twelve ordinary symbols. Their signed excesses are
divisible by eleven.

Within a replaced block, $B_1$ and $B_2$ each give ten closed trails,
since their excesses are divisible by ten. The first source trail needs
a different calculation. Its excess is 126, and a single lap sends an
insertion position to $p-126=p+4\pmod{10}$. This permutation has two
orbits,
\[
 (0,4,8,2,6),\qquad (1,5,9,3,7).
\]
Follow the transported paths through five successive laps of $B_0$,
starting at 0 and at 1. The two resulting trails are closed and together
use every transported row exactly once. Five laps have excess 630, so
each resulting trail has excess $630\cdot11/10=693$, divisible by eleven.
No row is duplicated by this procedure. The five laps follow five
different insertion positions, and the two orbits partition all ten.

A replaced block therefore gives $2+10+10+966=988$ trails. Its first
22 trails come from the three source trails; the other 966 are full
charts. Altogether the selection has
\[
 \num{10080}+48(988)=\num{57504}
\]
closed trails, all with signed excess divisible by eleven. Of these,
$\num{11136}$ come from the five source trails and $\num{46368}$ are
full charts. Completion gives eleven closed trails from each parent.

\subsection{The connector certificate}

For each of the $\num{46368}$ full charts, use the connector cycle of
its chart word. This meets all eleven of the chart's completed trails,
and it has the full-chart incidence required by
Proposition~\ref{prop-transport}.

The other $\num{11136}\cdot11=\num{122496}$ completed trails are met by
$\num{14774}$ additional connector cycles. These cycles have length eleven,
contain the completion symbol 13, and exclude the distinguished symbol 12.
The finite certificate gives one incidence witness for every pair
consisting of a parent trail and a completion position. Thus the total
number of connectors is
\[
 c_0=\num{46368}+\num{14774}=\num{61142}.
\]

An incidence witness is more than a label asserting that two objects
meet. It specifies a row occurrence in one of the five source trails,
the insertion positions that lead to the chosen parent, and a rotation
of one of the listed connector words. Prefix signed excesses determine
the position on entering the chosen row. The witness checks that the
head of the resulting completed row is exactly that rotation of the
connector word.

For the unchanged Halfar part, the occurrence passes through four
transports. For a replaced block it passes through one transport and
a relabeling. An occurrence descending from $B_0$ also specifies which
of the five laps is used. Before that lap, the prefix excess includes
126 for each preceding lap. The transport identities prove that these
instructions identify actual row occurrences in the stated parent trail.
The completion identities then prove the claimed endpoint equality.

The complete witness table is checked in Lean. Its positions fix the
parent and the completion position independently of the witness stored
there, so all $\num{122496}$ required pairs are covered. The circle
indices are checked against the literal list of $\num{14774}$ words.
This proves the second connector count and all the incidence hypotheses
of Proposition~\ref{prop-transport}.

\section{The finite and asymptotic bounds}\label{sec-proof}

The construction of Section~\ref{sec-base} satisfies
Proposition~\ref{prop-transport} with $n_0=14$. Its normalized costs are
\[
 \frac{Q_0}{11!}=\frac{889}{1800},\qquad
 \frac{c_0}{10!}=\gam.
\]
Their sum is
\[
 \frac{889}{1800}+\gam=\lam.
\]
Substituting these two costs in the proposition proves
Theorem~\ref{thm-finite}.

To deduce Theorem~\ref{thm-asymptotic}, put $m=n-2$ and choose an integer
$a$ with $2\le a\le m$, so that $\ell=m-a$. After division by
$(n-3)!=(m-1)!$, the two extra terms in the finite bound are at most
\begin{equation}\label{eq-error}
 \gam\frac{a-2}{m-1}
       +\frac{(m-a)m(m+1)}{(a+1)!}.
\end{equation}
Take $a=\lceil\sqrt m\rceil$. The first term tends to zero. For large
$a$, the last eight factors of $(a+1)!$ are each at least $a/2$, so the
second term is at most $2^8m^2(m+1)/a^8$, which also tends to zero.
This proves the theorem.

The formal finite statement uses the parameter $a=n-2-\ell$ and supplies
an actual covering word on an arbitrary finite alphabet. Its asymptotic
consequence is stated for positive rational $\varepsilon$; density of the
rationals gives the formulation for real $\varepsilon$ above. The five
source trails, the block coverage checks, and the entire incidence table
are verified by ordinary Lean kernel reduction. No step in this upper
proof uses \texttt{native\_decide}. The final declarations use only
\texttt{propext}, \texttt{Classical.choice}, and \texttt{Quot.sound}.
The programs that exported the finite data are not proof assumptions.
The Palomar submission selects Theorems~\ref{thm-finite}
and~\ref{thm-asymptotic}; the numerical corollaries are supporting material.
The public statements and their proofs are in
\path{Challenge.lean} and
\path{Solution.lean}.
The finite data and their checks are in \path{upper/}.

\section{Explicit words on ten through thirteen symbols}\label{sec-words}

The general theorem uses a uniform estimate for the cost of joining
components. For a particular alphabet, one can choose the cuts and the
order directly. Table~\ref{tab-words} gives the four explicit words
distributed with this paper.

\begin{table}[ht]
\centering
\begin{tabular}{rrr}
\toprule
$n$ & Word length & Length minus $\F(n)$\\
\midrule
10 & \num{4034857} & \num{2857}\\
11 & \num{43930548} & \num{22068}\\
12 & \num{522737007} & \num{189807}\\
13 & \num{6747802736} & \num{1863536}\\
\bottomrule
\end{tabular}
\caption{Explicit words distributed with this paper. Their coverage is
checked directly, separately from the Lean existence theorems.}
\label{tab-words}
\end{table}

The optimization can be separated into two decisions. Once the cyclic
pieces and their order are fixed, the choice of openings is a finite
shortest-path problem. Put every permitted opening of the $i$th piece
in layer $i$. The edge to an opening of the next piece has weight equal
to the number of new letters needed at that join. Dynamic programming
finds the best choices within this fixed order and the specified set of
openings. Searching over orders, moving groups of pieces, and changing
local images of a selection provide new inputs to this exact calculation.
Exactness for a fixed order does not imply optimality over all words.

At ten symbols the selection has 350 closed trails, obtained from our
earlier selection with 364 trails by a pair switch. The organization of
the mixed trails follows rumstd's contribution~\cite{rumstd}. We choose
the arrangement of groups of ordinary trails and mixed-trail runs, and
then their openings, by dynamic programming within the stated arrangement
model.

At eleven symbols the starting selection and arrangement are Halfar's
selection X. A block may be relabeled by rotations of its old symbols,
by exchanging its two added symbols, or by reversal with the corresponding
relabeling. These operations preserve its internal counts but change the
available ends. We vary these block images together with the arrangement
and then optimize the openings. The twelve-symbol word uses Halfar's
selection E and starts from his plan of length $\num{522737299}$. The
thirteen-symbol word uses his selection N and generator. In both cases
we optimize the order and cuts.

Every word was scanned directly for permutation coverage. The accompanying
checker ranks the length-$n$ factors with distinct symbols and verifies
that all $n!$ ranks occur. The repository contains the checker, the four
check records, and the exact words. These literal checks are separate
from the Lean proof of Theorems~\ref{thm-asymptotic} and~\ref{thm-finite}.

The files use the first $n$ characters of \word{0123456789ABC}, one byte
per symbol, without a final newline. The ten- and eleven-symbol words are
plain text; the two larger words are compressed without loss. A manifest
records their lengths and SHA-256 hashes, so the checked words can be
identified independently of their compressed representation. Instructions
for decompression and verification are in the repository's word guide.

The group discussion of October 2026 records further finite improvements.
Echols introduced compact recipes for fourteen-symbol words and a move that
reopens consecutive trails. Theo H. varied trail placement and opening phase,
and Halfar combined these moves with his own selections and searches~\cite{discussion}.
Halfar announced lengths $\num{522736969}$, $\num{6747802319}$, and
$\num{93904804425}$ for twelve, thirteen, and fourteen symbols, respectively.
The first two are shorter than the corresponding words distributed here;
the last improves the fourteen-symbol value in Table~\ref{tab-finite}.
Our eleven-symbol word has $30$ fewer letters than Halfar's announced
$\num{43930578}$. These comparisons keep the explicit words separate from
the general existence theorem. At eight symbols, Halfar's announced
$\num{46174}$ was followed by Theo H.'s $\num{46173}$, obtained from the
same selection with different openings and joins~\cite{eight}.

\begin{thebibliography}{9}

\bibitem{pantone}
J. Pantone,
\emph{Superpermutations with $(n-3)!$ coefficient $43/80$},
2026. Construction, exposition, and Lean formalization available at
\url{https://github.com/jaypantone/superperm-upper-43-80}.

\bibitem{halfar}
J. Halfar,
\emph{Superpermutations}, October 2026 source and construction notes.
The selections, block replacements, and coefficient $1771/3456$
used here are recorded at
\url{https://github.com/jhalfar/superpermutations/tree/cc5c763bd5854317b20dadb138990802656e6f03}.
The coefficient was announced on 6 October 2026 at
\url{https://groups.google.com/g/superpermutators/c/6ozhk3faHDA}.

\bibitem{rumstd}
rumstd,
\emph{Shorter explicit superpermutations}, contribution to Pantone's
construction, 2026,
\url{https://github.com/jaypantone/superperm-upper-43-80/pull/3}.

\bibitem{discussion}
W. Echols, Theo H., and J. Halfar,
\emph{Shorter words and compact recipes for superpermutations},
Superpermutators discussion, 17--20 October 2026,
\url{https://groups.google.com/g/superpermutators/c/iIgHJwUc9ww}.

\bibitem{eight}
J. Halfar and Theo H.,
\emph{An eight-symbol superpermutation of length 46,174 and an improvement to 46,173},
Superpermutators discussion, 20 October 2026,
\url{https://groups.google.com/g/superpermutators/c/a-oOgxxHHHg}.

\end{thebibliography}
\end{document}
